% ============================================================================
% 08_discussion.tex — Discussion
% ============================================================================

\section{Discussion}
\label{sec:discussion}

\subsection{International context}

The near-parity result is a product of the particular history England delivered. Table~\ref{tab:bis} places English house-price growth in international context using BIS residential property price statistics \parencite{bis_rpp}. The United Kingdom sits mid-table over 2005--2025: cumulative nominal growth of 95\% (3.3\% per year), far below the central European booms and well below the pace that would have made ownership financially dominant. 

\begin{table}[H]
  \centering
  \caption{Nominal house-price growth, selected countries, 2005Q1--2025Q4}
  \label{tab:bis}
  \begin{threeparttable}
  \begin{tabularx}{0.7\textwidth}{@{}Xrr@{}}
    \toprule
    \textbf{Country} & \textbf{Cumulative} & \textbf{Annualised} \\
    \midrule
    Hungary          & 356\% & 7.6\% \\
    Portugal         & 156\% & 4.6\% \\
    Netherlands      & 110\% & 3.6\% \\
    United Kingdom   & 95\%  & 3.3\% \\
    United States    & 89\%  & 3.1\% \\
    Germany          & 85\%  & 3.0\% \\
    Ireland          & 62\%  & 2.3\% \\
    France           & 60\%  & 2.3\% \\
    Spain            & 60\%  & 2.3\% \\
    Italy            & 15\%  & 0.7\% \\
    \bottomrule
  \end{tabularx}
  \begin{tablenotes}
    \small
    \item \textit{Source:} BIS residential property prices, nominal, selected series. The UK entry covers the whole United Kingdom; the representative English series of this paper grew 96\% (3.2\% per year) over the same span.
  \end{tablenotes}
  \end{threeparttable}
\end{table}

\subsection{Limitations}

Several limitations qualify the results. \emph{Aggregation.} The representative dwelling is a construct; households decide in a specific market, and the regional results of Section~\ref{sec:regional}, with buy-win shares from 38\% to 63\%, may differ from the aggregate result. Even within regions, composition differences between the owned and rented stocks mean the price-to-rent \emph{level} carries measurement error. \emph{Behavioural assumptions.} The renter invests 100\% of every monthly difference, immediately and for two decades, in a globally diversified portfolio. The saving and portfolio-behaviour literature \parencite{laibson1997,thaler2004} puts this well above typical conduct; the FTSE-only robustness check shows how costly even a diversification failure is, before any failure of discipline. \emph{Taxes.} The symmetric no-tax treatment is appropriate for a typical first-time buyer (no capital-gains tax on a main residence; ISA capacity covering the renter's flows for most of the sample) but understates the drag on very large taxable portfolios, which would shade outcomes towards buying. \emph{Inference.} The 199 cohorts overlap and share the sample's single realisation of asset returns: 258 months contain only 4.3 non-overlapping five-year windows, so cohort counts are counts of simulations rather than of independent observations. Win shares and correlations describe this history and are not tests on independent draws. \emph{Unpriced attributes.} Ownership insures against rent risk \parencite{sinai2005} and carries control rights; renting offers more flexibility and mobility \parencite{ehs2024}. None of these enters the terminal-wealth accounting. 

\subsection{Implications}

Three implications follow for households and for the framing of housing-policy debates. First, the popular certainty that buying is the financially correct choice is unwarranted, and not only at the five-year horizon. Across the 199 five-year cohorts the two strategies finish level: a median wealth ratio of 0.98, a median gap of \gbp{994} in the owner's favour and a mean gap of \gbp{1{,}744}. Lengthening the hold tilts the count towards buying, which wins 62\% of seven-year and 61\% of ten-year cohorts, but the median renter still finishes within 8--11\% of the owner, a gap of \gbp{14{,}454} at ten years on a dwelling worth \gbp{319{,}488}. Buying is ahead more often than not beyond five years, and it is not ahead by much.

Second, entry valuation and financing cost largely determine the cohort outcome without forecasting it, for the reasons Section~\ref{sec:predictors} sets out. The one regularity that survives runs in a single direction: renting won every cohort entering above a price-to-rent ratio of 24.4, and there is no matching floor.

Third, the comparison is conditional on conduct: the renter's side of the parity requires actually investing the cost difference, every month. A mortgage, besides providing crucial leverage, functions as a commitment device, since amortisation is forced saving. For households unlikely to match that discipline voluntarily, home ownership is likely the safer default even where the arithmetic is level, and that behavioural asymmetry is probably the strongest practical argument for buying.